\documentclass[10pt,a4paper]{amsart}
\usepackage[margin=19mm]{geometry}
\usepackage{amsmath,amssymb,mathtools}
\usepackage[scr=rsfs,cal=euler]{mathalfa}
\usepackage{xfrac}
\usepackage[unicode,hidelinks,bookmarks=false]{hyperref}
\newcommand{\ie}{i.e.\ }
\newcommand{\cf}{cf.\ }
\newcommand{\resp}{respectively\ }
\newcommand{\Subsection}{\subsection}
\providecommand{\nobreakdash}{\nobreak\hbox{-}}
\setlength{\emergencystretch}{2em}
\multlinegap0pt
\usepackage{amssymb}
\usepackage[shortlabels]{enumitem}
\newtheorem{assumption}{Assumption}
\newtheorem{theorem}{Theorem}
\newtheorem{lemma}{Lemma}
\DeclareMathOperator{\Ent}{Ent}
\DeclareMathOperator{\Law}{Law}
\newcommand{\Prob}{\mathcal{P}}
\newcommand{\R}{\mathbb{R}}
\newcommand{\Wass}{\mathcal{W}}
\title{A Mean-Field Theory of Transformers:\\ Well-Posedness of the Coupled Data--Parameter Dynamics\\ and Global Convergence of Training}
\author{Michael Herty, Hailiang Liu}
\date{Source: arXiv:2608.25055v1 (CC BY 4.0)}
\begin{document}
\maketitle

\noindent\emph{Source and licence.} A Mean-Field Theory of Transformers:\\ Well-Posedness of the Coupled Data--Parameter Dynamics\\ and Global Convergence of Training. Version arXiv:2608.25055v1 (CC BY 4.0), distributed by arXiv under the Creative Commons Attribution 4.0 International licence, http://creativecommons.org/licenses/by/4.0/, as recorded in the arXiv licence metadata for this exact version and shown on the abstract page https://arxiv.org/abs/2608.25055v1. Full licence notice: \texttt{licenses/arxiv-2608.25055-CC-BY-4.0.txt}.\par\smallskip

\noindent\textit{Editorial scope.} Counted unit: the article's theorem thm:ntk-conv (source0.tex lines 1962--1976). The article's complete original proof of it is delivered with it (source0.tex lines 1979--2159, one \texttt{proof} environment of 181 lines). No mathematics is written, repaired or re-derived: the statement and the proof are the article's own lines, in the article's own notation, and no retained line is substituted or re-flowed.  Retained uncounted as the source-local context the proof needs: lines 207--210; lines 225--228; lines 242--246; lines 248--251; lines 254--257; lines 299--303; lines 890--893; lines 1116--1119; lines 1121--1128; lines 1291--1299; lines 1734--1741; lines 1940--1942.  Nothing was omitted from any retained span, so the package carries no omission record.  No retained line contains a citation, so the wrapper carries no bibliography and no bibliography entry.  No retained line contains a figure environment, an image-inclusion command or drawing code, so no image is delivered and the package declares no asset dependency.  Editorial formatting only, in addition: an AMS \texttt{amsart} preamble that restates the article's own declarations and macros used by the retained lines, namely the environments \texttt{assumption}, \texttt{theorem}, \texttt{lemma} on separate counters, together with the standard packages those lines need.  The retained environments print in this document as Assumption 1, Theorem 1, Lemma 1 in order of appearance, whereas the article prints the same items with the numbering of its own full document.  The complete native source of the article as distributed in the arXiv e-print for this exact version is bundled unchanged as \texttt{sources/208539/source0.tex}.
\begin{equation}
K_\mu(x,y) = \frac{\exp\big(\tfrac1{\sqrt d}\langle Qx,Ky\rangle\big)}{\int \exp\big(\tfrac1{\sqrt d}\langle Qx,Ky'\rangle\big)\,d\mu(y')}.
\label{eq:kernel}
\end{equation}

\begin{equation}
V_{\rho,\mu}(x) := \int_\Theta A(x,\mu,\theta)\,\rho(d\theta),
\label{eq:V}
\end{equation}

\begin{equation}
\partial_s\rho - \nabla_\theta\cdot \left(\rho  \nabla_\theta\frac{\delta\mathcal L}{\delta\rho}(\theta,\rho) \right) = 0, 
%\qquad u(\theta,s) = -\nabla_\theta\frac{\delta\mathcal L}{\delta\rho}(\theta,\rho_s),
\label{eq:training-transport}
\end{equation}

\begin{equation}
\partial_s\rho = \nabla_\theta\cdot\Big(\rho\,\nabla_\theta\frac{\delta\mathcal L}{\delta\rho}\Big) + \beta\Delta_\theta\rho.
\label{eq:mfld}
\end{equation}

\begin{equation}
\mathcal F_\beta(\rho) := \mathcal L(\rho) + \beta\,\Ent(\rho) , \qquad \Ent(\rho):=\int \rho\log\rho\, d\theta.
\label{eq:Fbeta}
\end{equation}

\begin{assumption}[Compact parameter domain for the forward pass]
\label{ass:compact-theta}
$\Theta\subset\R^p$ is compact, with $\|\theta\|\le D_\Theta$ for all $\theta=(Q,K,V)\in\Theta$ (equivalently, we work with any fixed, or randomly initialized but almost-surely bounded, family of heads. 
%this is automatic under the usual compactly-supported or sub-Gaussian initializations).
\end{assumption}

\begin{equation}
\frac{\delta\mathcal L}{\delta\rho}(\theta) = \int_0^T\!\!\int_{\R^d}  \nabla_x p(t,x)\cdot A(x,\mu_t,\theta)\,d\mu_t(x)\,dt.
\label{eq:grad-rho}
\end{equation}

\begin{assumption}[Confinement]
\label{ass:confine}
$\mathcal L$ in \eqref{eq:mfld} is replaced by $\mathcal L_\lambda(\rho):=\mathcal L(\rho) + \tfrac\lambda2\int|\theta|^2\,d\rho(\theta)$ for some $\lambda>0$ (an $\ell_2$/weight-decay regularizer), and $\theta\mapsto \nabla_\theta A(x,\mu,\theta)$ (hence $\theta\mapsto\delta\mathcal L/\delta\rho(\theta)$, by \eqref{eq:grad-rho}) is $L_A$-Lipschitz and $L_A$-bounded in $\theta$, uniformly over $\mu$ supported in $B_{R(T)}$; this holds e.g.\ whenever $A(x,\mu,\theta)$ is twice continuously differentiable in $\theta$ with bounded second derivative on bounded sets, as is the case for the softmax-attention kernel \eqref{eq:kernel}--\eqref{eq:V}.
\end{assumption}

\begin{theorem}[Existence and uniqueness of the mean-field Langevin training dynamics]
\label{thm:mfld-wp}
Under Assumptions~\ref{ass:compact-theta}--\ref{ass:confine}, for every $\rho_0\in\Prob_2(\R^p)$ with finite entropy, Eq.~\eqref{eq:mfld} (with $\mathcal L$ replaced by $\mathcal L_\lambda$) has a unique weak solution $\rho\in C([0,\infty);\Prob_2(\R^p))$ with $\sup_{s\ge0}\big(\mathbb E_{\rho_s}[|\theta|^2] + \Ent(\rho_s)\big)<\infty$. Equivalently, $\rho_s=\Law(\Theta_s)$ where $(\Theta_s)_{s\ge0}$ solves the McKean--Vlasov SDE
\begin{equation}
d\Theta_s = -\nabla_\theta\frac{\delta\mathcal L_\lambda}{\delta\rho}(\Theta_s,\rho_s)\,ds + \sqrt{2\beta}\,dW_s, \qquad \Law(\Theta_0)=\rho_0.
\label{eq:mfld-sde}
\end{equation}
\end{theorem}

\begin{lemma}[Dissipation identity]
\label{lem:dissipation}
Along any solution of \eqref{eq:mfld} given by Theorem~\ref{thm:mfld-wp},
\begin{equation}
\frac{d}{ds}\mathcal F_\beta(\rho_s) = -\int_\Theta \Big|\nabla_\theta\frac{\delta\mathcal F_\beta}{\delta\rho}(\theta,\rho_s)\Big|^2\,d\rho_s(\theta) \le 0,
\label{eq:dissipation}
\end{equation}
where $\mathcal F_\beta$ is the regularized objective \eqref{eq:Fbeta}  with $\mathcal L$ or  $\mathcal L_\lambda$, respectively.
\end{lemma}

\begin{assumption}[NTK non-degeneracy at initialization]
\label{ass:ntk}
Assume that the NTK operator is strictly positive on the mean-zero subspace at initialization. That is, $\lambda_0 := \lambda_{\min}(G(\rho_0)) > 0$. 
 Assume that $\rho\mapsto G(\rho)$ is $L_G$-Lipschitz on $\Prob(\Theta)$ in the sense that 
 $$
 \|G(\rho)-G(\rho')\|_{\mathrm{op}}\le L_G\,\Wass_2(\rho,\rho').
 $$
\end{assumption}

\begin{equation}\label{phir} 
\phi_\rho(\theta) := \int_0^T \int \nabla_x q_\rho(t, x)\cdot A(x, \mu_t^\rho, \theta)d\mu_t^\rho dt.
\end{equation} 

\begin{theorem}[Local linear convergence via NTK positivity]
\label{thm:ntk-conv} Consider the loss function of form 
$$
J(\mu_T)=\frac{1}{2} \left|y- \int x d\mu_T(x) \right|^2. 
$$
Let Assumptions~\ref{ass:compact-theta}--\ref{ass:confine}, \ref{ass:ntk} hold. Assume a pure Wasserstein gradient flow \eqref{eq:training-transport} (i.e.\ $\beta=0$ in \eqref{eq:mfld}) and assume that  the initial loss satisfies
\begin{equation}\label{ini}
\mathcal L(\rho_0) - \min_\rho \mathcal L(\rho) \le \frac{ \lambda_0^4}{16  L_K (\lambda_0+2\Lambda L_K)},
\end{equation}
then the gradient flow $\rho_s$ of Theorem~\ref{thm:mfld-wp} with $\beta=0$ satisfies, for all $s\ge0$,
\begin{equation}
\mathcal L(\rho_s) - \min_\rho\mathcal L(\rho) \le e^{-\lambda_0 s}\,\big(\mathcal L(\rho_0)-\min_\rho\mathcal L(\rho)\big).
\end{equation}
Hence, the training dynamics converges at a linear rate to a global minimizer of $\mathcal L$.
\end{theorem}

\begin{proof} The proof consists of three steps. \\ 
{\bf Step 1: Preservation of NTK non-degeneracy.}\\
By Assumption~\ref{ass:ntk} and the Lipschitz continuity of $K_\rho$ we have 
\[
\lambda_{\min}(K_{\rho_s})
\ge
\lambda_{\min}(K_{\rho_0})
-
L_K
\Wass_2(\rho_s,\rho_0) \geq \lambda_0 - L_K
\Wass_2(\rho_s,\rho_0). 
\]
Hence, as long as 
\[
\Wass_2(\rho_s,\rho_0)
\le
\frac{\lambda_0}{2L_K},
\]
we have
\[
\lambda_{\min}(K_{\rho_s})
\ge
\frac{\lambda_0}{2}.
\]
Thus, the NTK remains uniformly coercive as long as the training trajectory remains inside a sufficiently small Wasserstein neighborhood of its initialization. Moreover,
$$
\lambda_{\max}(K_{\rho_s}) \leq \Lambda +L_K
\Wass_2(\rho_s,\rho_0) \leq \Lambda + \frac{\lambda_0}{2L_K}. 
$$

{\bf Step 2: Energy dissipation and the Polyak–Łojasiewicz inequality.} \\ 

By the dissipation identity given in Lemma~\ref{lem:dissipation}  with $\beta=0$, 
\begin{equation}
\frac{d}{ds}\mathcal L(\rho_s)
=
-
\int 
\bigl|
\nabla_\theta\Psi_{\rho_s}(\theta)
\bigr|^2
\,d\rho_s(\theta).
\label{eq:pl-step1}
\end{equation}
We first show that NTK non-degeneracy  implies the Wasserstein Polyak–Łojasiewicz 
inequality 
\[
\left\|\nabla_\theta \Psi_\rho\right\|_{L^2(\rho)}^2
\ge \lambda_0
\bigl(
\mathcal L(\rho)
-
\mathcal L^\star
\bigr),
\]
where
\(
\mathcal L^\star
=
\min_\rho\mathcal L(\rho).
\)
%Then $J(\mu_T)=\frac{1}{2}(y-m(\mu_T))^2, \quad m(\mu):= \int x d\mu(x)$. 
Let $m_\rho$ denote the first momentum of the terminal state $m(\mu_T^\rho)$, with 
$$
m(\mu):=\int x d\mu(x) \in \mathbb{R}^d, 
$$
and define the  vector residual 
$r_\rho:=m_\rho - y$. We consider the terminal loss  
$$
L(\rho)=J(\mu_T^\rho)=\frac{1}{2}|r_\rho|^2. 
$$
A direct calculation gives  
$$
\frac{\delta J}{\delta \mu_T}(x)= r_\rho \cdot x. 
$$
Hence the terminal condition for the adjoint equation is 
 $$
 p(T, x)=\frac{\delta J}{\delta \mu_T}(x)= r_\rho \cdot x.
 $$ 
 For fixed $\rho$, the adjoint equation is linear in $p$, write  
$$
p(t, x)=r_\rho \cdot q_\rho(t, x), 
$$
where $q_\rho(t, x)\in \mathbb{R}^d$ is vector-valued and solves the same adjoint equation with normalized terminal condition $q_\rho(T, x)=x$. Consequently,  
$$
\nabla_x p(t, x)=( \nabla_x q_\rho(t, x))^\top r_\rho. 
$$
Substituting this representation into the formula for the first variation of $\mathcal L$ gives
$$
\Psi_\rho(\theta)=\frac{\delta \mathcal L}{\delta \rho}(\theta; \rho)= (\phi_\rho(\theta))^\top r_\rho,
$$
where 
$
\phi_\rho(\theta)
$
is given in (\ref{phir}).  Since $r_\rho$ is independent of $\theta$,
$$
\nabla_\theta \Psi_\rho(\theta)=(\nabla_\theta \phi_\rho(\theta))^\top r_\rho. 
$$
Thus 
$$
\|\nabla_\theta \Psi_{\rho_s}\|^2_{L^2(\rho_s)}= r_{\rho_s}^\top K_{\rho_s} r_{\rho_s}  \geq \frac{\lambda_0}{2} |r_{\rho_s}|^2. 
%|r_\rho|^2 \int |\nabla_\theta \phi_\rho(\theta)|^2 d\rho(\theta).  
$$
If the minimum loss is attained at zero, so that $\mathcal L^*=0$, then 
$$
|r_{\rho_s}|^2= 2( \mathcal L(\rho_s)-\mathcal L^*). 
$$
Therefore 
$$
\|\nabla_\theta \Psi_{\rho_s}\|^2_{L^2(\rho_s)} \geq  \lambda_0  (\mathcal L(\rho)-\mathcal L^* ). 
$$
This establishes the desired Wasserstein PL inequality. Combining this inequality with the dissipation identity (\ref{eq:pl-step1}), we obtain  
\[
\frac{d}{ds}
\bigl(
\mathcal L(\rho_s)-\mathcal L^\star
\bigr)
\le
-
\lambda_0
\bigl(
\mathcal L(\rho_s)-\mathcal L^\star
\bigr),
\]
and Gr\"onwall's inequality immediately implies
\[
\mathcal L(\rho_s)-\mathcal L^\star
\le
e^{-\lambda_0 s}
\bigl(
\mathcal L(\rho_0)-\mathcal L^\star
\bigr).
\]
{\bf Step 3: $\rho_s$ staying in the neighborhood.}\\
It remains to verify that the trajectory never leaves the neighborhood
$\Wass_2(\rho_s,\rho_0)\le \lambda_0/(2L_k)$ on which the above argument is valid.
Since the metric derivative of the Wasserstein gradient flow satisfies
\[
|\partial_s  \rho_s|_{\Wass_2}
=
\|
\nabla_\theta\Psi_{\rho_s}
\|_{L^2(\rho_s)},
\]
the total displacement satisfies
\[
\Wass_2(\rho_0,\rho_s)
\le
\int_0^s
\|
\nabla_\theta\Psi_{\rho_\tau}
\|_{L^2(\rho_\tau)}
\,d\tau.
\]
Using Assumption ~\ref{ass:ntk} again,  we have the upper bound 
$$
\|\nabla_\theta \Psi_{\rho_s}\|^2_{L^2(\rho_s)}= r_{\rho_s}^\top K_{\rho_s} r_{\rho_s} \leq (\Lambda +\lambda_0/(2L_K))  |r_{\rho_s}|^2 
=(2 \Lambda+\lambda_0/L_K) ( \mathcal L(\rho_s)-\mathcal L^*).
$$
Combining this with the exponential decay established above 
$$
 \mathcal L(\rho_s)-\mathcal L^*)  \leq  e^{-\lambda_0 s} ( \mathcal L(\rho_0)-\mathcal L^*),
$$
we obtain %Using the exponential decay established \textbf{}above together with the energy
%dissipation identity and Cauchy--Schwarz yields
$$
\|\nabla_\theta \Psi_{\rho_s}\|_{L^2(\rho_s)}
\leq \sqrt{2\Lambda +\lambda_0/L_K} e^{-\lambda_0 s/2} ( \mathcal L(\rho_0)-\mathcal L^*)^{1/2}. 
$$
Consequently,  
\[
\sup_{s\ge0}
\Wass_2(\rho_0,\rho_s)
\le
\frac{2\sqrt{2\Lambda +\lambda_0/L_K} }{\lambda_0}
(\mathcal L(\rho_0)-\mathcal L^\star)^{1/2}.
\]
Hence, if the initial excess loss is sufficiently small as shown in (\ref{ini}),  the entire trajectory remains inside the neighborhood $\Wass_2(\rho_s,\rho_0)\le \lambda_0/(2L_k)$, 
where  the NTK non-degeneracy and upper-bound assumptions remain valid.
\end{proof}

\end{document}
